% Einheit 3 von 3 – Kugel (bank/pyramide-kegel-kugel/e3.jsonl, zusammenbau v0.1)
%% TODO Zweigzeile Teil 1: was hier gelernt wird (Satz in Schülersprache aus dem Einheitstitel) – Einheit 3
\einheitenkopf[e3]{Einheit 3 von 3 · Kugel}
\zweigzeile{ab Kl. 8, je nach Buch bis Kl. 10 (am Gymnasium ab Kl. 9) · P10}
\setcounter{aufgabe}{20}

%% TODO Titel: Ich-kann-Satz fehlt in der Bank (Platzhalter: Kettenname) – Nr. 21
%% TODO Anweisung: ein Satz über der ersten Teilaufgabe fehlt in der Bank – Nr. 21
\begin{aufgabe}{Kugel}
\begin{teile}
\teil Die Strecke mit ? an der Kugel – Radius oder Durchmesser? Kreuze an.\\ \kreuz{Radius}\\ \kreuz{Durchmesser}
\end{teile}
%% TODO Darstellung für schwach fehlt in der Bank (pyramide-kegel-kugel-e3-k1-s1-v1; Abschnitt „Für schwache Schüler“)
\swz{Kugel: Radius $3$ cm – Volumen? Runde auf eine Stelle.}{}
%% TODO Darstellung für schwach fehlt in der Bank (pyramide-kegel-kugel-e3-k1-s1-v2; Abschnitt „Für schwache Schüler“)
\swz{Kugel: Radius $4$ cm – Volumen? Runde auf eine Stelle.}{}
%% TODO Darstellung für schwach fehlt in der Bank (pyramide-kegel-kugel-e3-k1-s1-v3; Abschnitt „Für schwache Schüler“)
\swz{Kugel: Radius $2$ dm – Volumen? Runde auf eine Stelle.}{}
%% TODO Darstellung für schwach fehlt in der Bank (pyramide-kegel-kugel-e3-k1-s1-v4; Abschnitt „Für schwache Schüler“)
\swz{Kugel: Radius $7$ cm – Volumen? Runde auf eine Stelle.}{}
%% TODO Darstellung für schwach fehlt in der Bank (pyramide-kegel-kugel-e3-k1-s1-v5; Abschnitt „Für schwache Schüler“)
\swz{Kugel: Radius $9$ m – Volumen? Runde auf eine Stelle.}{}
\swfrage{Was bleibt gleich, was ändert sich?}
%% TODO Darstellung für schwach fehlt in der Bank (pyramide-kegel-kugel-e3-k1-s2-v1; Abschnitt „Für schwache Schüler“)
\swz{Kugel: Durchmesser $16$ cm – Volumen? Runde auf eine Stelle.}{}
%% TODO Darstellung für schwach fehlt in der Bank (pyramide-kegel-kugel-e3-k1-s3-v1; Abschnitt „Für schwache Schüler“)
\swz{Kugel: Radius $2{,}5$ cm – Volumen? Runde auf eine Stelle.}{}
%% TODO Darstellung für schwach fehlt in der Bank (pyramide-kegel-kugel-e3-k1-s4-v1; Abschnitt „Für schwache Schüler“)
\swz{Ein Wasserball hat den Radius $15$ cm. Wie viel Liter Luft sind darin?}{}
%% TODO Darstellung für schwach fehlt in der Bank (pyramide-kegel-kugel-e3-k1-s5-v1; Abschnitt „Für schwache Schüler“)
\swz{Kugel: Radius $11$ cm – Oberfläche? Runde auf eine Stelle.}{}
\swa{Skizziere eine Kugel mit dem Radius $2$ cm: Kreis, Äquator als Ellipse, Radius gestrichelt und beschriftet.}{\rechenplatz[halb]{4}}
%% TODO Darstellung für schwach fehlt in der Bank (pyramide-kegel-kugel-e3-k1-s7-v1; Abschnitt „Für schwache Schüler“)
\swz{Halbkugel: Radius $4$ cm – Volumen und Oberfläche? Runde auf eine Stelle.}{}
%% TODO Darstellung für schwach fehlt in der Bank (pyramide-kegel-kugel-e3-k1-s8-v1; Abschnitt „Für schwache Schüler“)
\swz{Eine Stahlkugel hat den Radius $2$ cm, Stahl hat die Dichte $7{,}9$ g/cm³. Wie schwer ist sie? Runde auf ganze Gramm.}{}
%% TODO Darstellung für schwach fehlt in der Bank (pyramide-kegel-kugel-e3-k1-s9-v1; Abschnitt „Für schwache Schüler“)
\swz{Kugel: Oberfläche $200$ cm² – Radius? Runde auf zwei Stellen.}{}
%% TODO Darstellung für schwach fehlt in der Bank (pyramide-kegel-kugel-e3-k1-s10-v1; Abschnitt „Für schwache Schüler“)
\swz{Ein Ball mit dem Radius $14$ cm soll genau in eine würfelförmige Schachtel passen. Wie lang ist die Kante? Wie viel cm³ Luft bleiben in der Schachtel?}{}
%% TODO Darstellung für schwach fehlt in der Bank (pyramide-kegel-kugel-e3-k1-s11-v1; Abschnitt „Für schwache Schüler“)
\swz{Ein Silo besteht aus einem Zylinder (Radius $3$ m, Höhe $9$ m) mit einem Halbkugeldach. Wie viel m³ fasst es insgesamt?}{}
%% TODO Darstellung für schwach fehlt in der Bank (pyramide-kegel-kugel-e3-k1-s12-v1; Abschnitt „Für schwache Schüler“)
\swz[3]{Kugel A hat den Radius $2$ cm, Kugel B den Radius $4$ cm. Berechne beide Volumen und zeige, dass B achtmal so viel Volumen hat.}{}
\swz[3]{Eine Bowlingkugel hat den Durchmesser $21{,}6$ cm. Berechne ihr Volumen mit der Formel aus der Formelsammlung \hfill (P10 2016 OS).}{}
\end{aufgabe}

%% TODO Titel: Ich-kann-Satz fehlt in der Bank (Platzhalter: Kettenname) – Nr. 22
%% TODO Anweisung: ein Satz über der ersten Teilaufgabe fehlt in der Bank – Nr. 22
\begin{aufgabe}{Kugel – Fehler finden, Begründen, Darstellungswechsel, Anwendung}
\swz[3]{Eine Kugel hat den Durchmesser $18$ cm. Nora rechnet das Volumen: \rechnung{V &= \tfrac{4}{3} \cdot \pi \cdot 18^3 \\ V &\approx 24\,429{,}0 \text{ cm}^3} Wo ist der Fehler?}{}
\swfrage{Warum ist die Oberfläche einer Halbkugel mehr als halb so groß wie die Oberfläche der ganzen Kugel? Begründe.}
\swa{Ein Silo besteht aus einem Zylinder mit dem Radius $2{,}5$ m und der Höhe $8$ m, oben sitzt ein Halbkugeldach. Skizziere das Silo und beschrifte Radius, Höhe des Zylinders und Gesamthöhe.}{\rechenplatz[halb]{4}}
\swz[3]{Die Kuppel eines Turms ist eine Halbkugel mit dem Radius $2{,}8$ m. Sie wird außen zweimal gestrichen; ein Eimer Farbe reicht für $15$ m². Wie viele Eimer braucht man?}{}
\end{aufgabe}

\uebersichtskasten{Volumen: vier Drittel mal π mal Radius hoch drei – der Radius steht dreimal in der Formel, weil der Raum drei Richtungen hat. \\ d = 10 cm, r = 5 cm: V = 4/3 · π · 5³ = 4/3 · π · 125 ≈ 523,6 cm³ ≈ 0,52 l \\ Oberfläche: vier mal π mal Radius hoch zwei – viermal der Kreis mit demselben Radius. \quad O = 4 · π · 5² ≈ 314,2 cm² \\ Halbkugel: Volumen halb; Oberfläche halb plus der Schnittkreis. \quad V ≈ 261,8 cm³ \quad O = 2 · π · 5² + π · 5² ≈ 235,6 cm² \\ Rückwärts: r = √(O : (4 · π)); aus dem Volumen mit der dritten Wurzel r = ³√(3 · V : (4 · π)).}
